\chapter{\ifproject%
\ifenglish Project Structure and Methodology\else โครงสร้างและขั้นตอนการทำงาน\fi
\else%
\ifenglish Project Structure\else โครงสร้างของโครงงาน\fi
\fi
}

ในบทนี้จะกล่าวถึงหลักการ และการออกแบบระบบไปจนถึงขั้นตอนการออกแบบจากความต้องการของผู้ใช้งาน

\makeatletter

% \renewcommand\section{\@startsection {section}{1}{\z@}%
%                                    {13.5ex \@plus -1ex \@minus -.2ex}%
%                                    {2.3ex \@plus.2ex}%
%                                    {\normalfont\large\bfseries}}

\makeatother
%\vspace{2ex}
% \titleformat{\section}{\normalfont\bfseries}{\thesection}{1em}{}
% \titlespacing*{\section}{0pt}{10ex}{0pt}

\section{การติดต่อและคุยงานเพื่อสรุปความต้องการของสำนักทะเบียน}

% \begin{figure}
% \begin{center}
% % \includegraphics{800px-Briny_Beach.jpg}
% \end{center}
% % \caption[Poem]{The Walrus and the Carpenter}
% \label{fig:walrus}
% \end{figure}

\section{เงื่อนไขของการวางร่างปฏิทินการศึกษา}
  เนื่องจากจุดประสงค์ของโครงงานนี้คือต้องการพัฒนาเว็บไซต์ให้แก่สำนักทะเบียน
  จึงจะต้องเริ่มจากการพูดคุยกับบุคลากรของสำนักทะเบียนเพื่อให้ได้ความต้องการที่แท้จริงของโครงงาน โดยในปฏิทินจะมีเงื่อนไขต่างๆ อันสรุปได้ดังนี้

\subsection{ภาคการศึกษาที่ 1}
 คาบเรียนของแต่ละวัน มีจำนวนดังนี้ วันจันทร์อย่างน้อย 14 คาบ วันอังคารอย่างน้อย 14 คาบ วันพุธอ่างน้อย 14 คาบ วันพฤหัสบดี 13 คาบ วันศุกร์ 12 คาบ(ประมาณ)\CIreply{is this from just one academic year, or from a requirement?}
\begin{itemize}
  \item วันเปิดภาคเรียนมักจะอยู่ในเดือนมิถุนายน
  \item รูปแบบวันจันทร์ พฤหัสบดี มีระยะเวลาการเรียนการสอนตลอดภาคการศึกษา จำนวน 42 ชั่วโมง
  \item รูปแบบวันอังคาร ศุกร์ มีระยะเวลาการเรียนการสอนตลอดภาคการศึกษา จำนวน 40 ชั่วโมง 30 นาที\CIreply{why?}
  \item วันลงทะเบียนเรียนล่วงหน้า สัปดาห์แรกของเดือนก่อนหน้าที่จะเปิดภาคการศึกษา
  \item ประกาศผลการลงทะเบียนเรียนล่วงหน้าหลังจากลงทะเบียนล่วงหน้าประมาณ 10 วัน
  \item วันลงทะเบียนเรียนมักจะเริ่มก่อนวันเปิดภาคเรียนวันแรก 1--2 วัน
  \item วันลงทะเบียนมีระยะเวลา 8 วัน (เป็นเวลาที่ตรงกับวันเรียนด้วย)
  \item วันถอนกระบวนวิชาโดยไม่ได้รับอักษรลำดับขั้น W จะเริ่มจากวันลงทะเบียนเรียน \CI{ถึงวันที่ 1 ของเดือนถัดไป}{really?}
  \item วันที่อาจารย์ที่ปรึกษาให้ความเห็นชอบการลงทะเบียนเรียนของนักศึกษาทางระบบ Internet ระยะเวลาเริ่มจาก 1 สัปดาห์หลังจากเปิดภาคเรียน\CI{จนถึงวันที่ 1 ของเดือนถัดไป}{really?}
  \item หลังจากหมดวันถอนกระบวนวิชาโดยไม่ได้รับอักษรลำดับขั้น W จะเริ่มนับวันถอนกระบวนวิชาโดยได้รับอักษรลำดับขั้น W
  \item หลังจากเปิดภาคเรียน เมื่อเรียนครบ 8 สัปดาห์จะเริ่มสอบกลางภาค
  \item สอบกลางภาคสัปดาห์ที่ 9 ของภาคเรียน มีระยะเวลาทั้งหมด 7 วัน\CIreply{what's different from last item?}
  \item หลังจากวันสุดท้ายของการสอบกลางภาค เริ่มเรียนครึ่งภาคเรียนหลัง (สัปดาห์ที่ 10)
  \item ครึ่งภาคเรียนหลังเวลาเรียนจะมีระยะเวลา 7 สัปดาห์ \CI{เริ่ม}{?}สอบกลางภาค
  \item วันสอบปลายภาคมีระยะเวลา 2 สัปดาห์
  \item หลังสอบปลายภาคจะทำการหยุดเรียนจนจบสัปดาห์นั้นและหยุดเรียนเพิ่มอีก 1 สัปดาห์ และเปิดภาคเรียนที่2วนวันจันทร์ของสัปดาห์ถัดมา
  \item วันไห้ครู ไม่มีการเรียนการสอน
\end{itemize}
  \CIreply{convert to itemize, as shown in the first few items above} 

\subsection{ภาคการศึกษาที่ 2}
\begin{itemize}
  \item คาบเรียนของแต่ละวัน มีจำนวนดังนี้ วันจันทร์ 12 คาบ วันอังคาร 15 คาบ วันพุธ 15 คาบ วันพฤหัสบดี 15 คาบ วันศุกร์ 14 คาบ
  \item รูปแบบวันจันทร์ พฤหัสบดี มีระยะเวลาการเรียนการสอนตลอดภาคการศึกษา จำนวน 40 ชั่วโมง 30 นาที
  \item รูปแบบวันอังคาร ศุกร์ มีระยะเวลาการเรียนการสอนตลอดภาคการศึกษา จำนวน 43 ชั่วโมง 30 นาที
  \item วันลงทะเบียนเรียนล่วงหน้า \CI{ก่อนเปิดภาคเรียน 25 วัน}{really?}
  \item ประกาศผลการลงทะเบียนล่วงหน้า \CI{1 สัปดาห์}{from when?}
  \item เปิดภาคการศึกษาที่ 2 (ระยะเวลา 18 สัปดาห์เหมือนภาคการศึกษาที่ 1)
  \item วันลงทะเบียนเรียนมักจะเริ่มก่อนวันเปิดภาคเรียนวันแรก 1--2 วัน
  \item วันลงทะเบียนมีระยะเวลา8วัน (เป็นเวลาที่ตรงกับวันเรียนด้วย)
  \item หลังจากสอบปลายภาคภาคเรียนที่2 จะหยุดเรียน3 สัปดาห์
  \item หลังจากหยุดเรียน3สัปดาห์ จะเริ่มเปิดภาคเรียนฤดูร้อน
\end{itemize}

\subsection{ภาคการศึกษาฤดูร้อน}
\begin{itemize}
  \item ภาคฤดูร้อน มีรูปแบบการเรียนการสอนตั้งแต่วันจันทร์--ศุกร์ ระยะเวลาที่ใช้สำหรับการเรียนการสอนแต่ละกระบวนวิชามีเวลาเรียนจำนวน 29 คาบ เท่ากับ 43 ชั่วโมง 30 นาที
  \item ลงทะเบียนล่วงหน้าของภาคฤดูร้อนเริ่มหลังจากปิดภาคเรียนที่ 2 ไปแล้ว 1 สัปดาห์(ระยะเวลา 4 วัน)
  \item ประกาศผลการลงทะเบียนล่วงหน้าหลังจากปิดการลงทะเบียน 2 วัน
  \item วันเพิ่ม-ถอนกระบวนวิชา/ลงทะเบียน สำหรับนักศึกษาทุกระดับ ก่อนวันเปิดภาคเรียน 1 สัปดาห์
  \item ภาคฤดูร้อนระยะเวลาเรียน 6 สัปดาห์
  \item สัปดาห์ที่ 7 ของภาคฤดูร้อนจะเป็นสอบปลายภาคฤดูร้อน
  \item หลังจากสอบปลายภาคฤดูร้อนแล้วจะปิดปีการศึกษา
\end{itemize}

\CI{หลังจากที่ได้เงื่อนไขทั้งหมดครบแล้ว จะนำเงื่อนไขเหล่านี้มาแยกออกจากระบบที่ตอบสนอง
และเพิ่มความสะดวกสบายของผู้ใช้}{what's the point of this paragraph?}

\section{การทำงานของโปรแกรม}
จากการสรุปความต้องการของสำนักทะเบียนผ่านทางบุคลากร จึงสรุปออกมาเป็น user flow หรือสิ่งที่แสดงเส้นทางของผู้ใช้แอพพลิเคชั่นได้ดังรูปที่~\ref{fig:user-flow-new}
\CIreply{this figure doesn't outline the entire user flow; did you mean to say something different?}

\begin{figure}[h]
\centering
\includegraphics[width=1\textwidth]{pic3.1.jpg}
\caption{การทำงานเมื่อผู้ใช้ต้องการสร้างแบบปฏิทินใหม่}
\label{fig:user-flow-new}
\end{figure}

โดยผู้ใช้จะเริ่มจากการเข้าสู่ระบบโดยใช้ CMU \CI{QAuth}{?} หลังจากนั้น คลิกที่สร้างดราฟใหม่ 
หลังจากนั้นเว็บไซต์จะต้องการทราบวันแรกของการเปิดภาคเรียนเพื่อนำไปสร้างปฏิทินการศึกษา
โดยหน้า Document จะเป็นหน้าที่ใช้จัดการกับร่างปฏิทินทั้งหมดที่ผู้ใช้ได้สร้างไว้
หลังจากที่ผู้ใช้ได้คลิกสร้างปฏิทินขึ้นมาใหม่ ระบบจะต้องการให้ผู้ใช้กรอกข้อมูลของวันเปิดเทอม
ของปีการศึกษานั้น หลังจากนั้นระบบจะทำการสร้างร่างปฏิทินการศึกษาแบบอัตโนมัติ
เพื่อทำให้ง่ายต่อการแก้ไข ไม่เกิดความยุ่งยากในการต้องมาเพิ่มกิจกรรมทีละวันกิจกรรม
โดยกิจกรรมที่นำไปใส่ลงในปฏิทินแบบอัติโนมัตินั้นจะได้มาจากการคำนวณ
วันที่อยู่ห่างจากวันเปิดเทอมตามเงื่อนไขของปฏิทิน
 วางรูป Dupilicate flow\
 รูปที่3.2 การทำงานเมื่อผู้ใช้ต้องการทำซ้ำปฏิทินเดิม

\section{โครงสร้างของโครงงาน}

\begin{figure}[h]
\centering
\includegraphics[width=1\textwidth]{pic3.2.jpg}
\caption{การทำงานเมื่อผู้ใช้ต้องการสร้างแบบปฏิทินใหม่}
\end{figure}

จากรูปที่~\CI{3.2}{use reference} ผู้ใช้ต้องการจะทำซ้ำ หลังจากที่ผู้ใช้อยู่ในหน้า Document และคลิกที่ทำซ้ำหรือ Duplicate
เว็บไซต์จะต้องการให้ผู้ใช้กรอกชื่อของปฏิทินที่จะสร้างใหม่ที่ทำซ้ำมาจากปฏิทินเดิม
และปีที่ต้องการเปลี่ยนใหม่ หากมีกิจกรรมของปฏิทินเดิมที่คล้ายคลึงกับปฏิทินของปีถัดไป 
ผู้ใช้สามารถทำซ้ำปฏิทินเดิมแล้วเปลี่ยนเป็นปีถัดไปได้เลย
\CIreply{what's different from the previous flow?}

\section{โครงสร้างโปรแกรม}
ในส่วนของ client จะใช้ภาษา React.JS ในการสร้างเว็บไซต์ แพลตฟอร์มนี้จะใช้กับ
คอมพิวเตอร์ โดยมี Node.JS ในส่วน backend และใช้ API ในการรับส่งกับฐานข้อมูล
และฐานข้อมูล MongoDB
